\section{Contrastive Learning：Agreement 是起点，不是答案}

Islands 的表面目标是“同一 entity 的不同位置得到相同向量”，这与 contrastive learning 很像。但普通 image-level objective 只知道两个 crop 来自同一张图，未必知道它们属于同一 scene、同一 object 还是不同 objects。GLOM 借用 agreement 与 anti-collapse 思想，却必须把它改造成 level-specific、local、pose-aware 的动态过程。

\lecturefigure{slide-15-contrastive-learning-intro.jpg}{Contrastive visual representation learning：用同图不同 patch 的 agreement 学表示。}{Stanford Online 官方视频 00:22:02--00:24:15。}

\subsection{SimCLR 的正样本、增强与 collapse}

上一节只说“agreement 可以学表示”，还没有说明如何避免所有向量一起相等。本节用 SimCLR 把这个缺口补齐：先明确 positive pair 的生成方式，再观察 naïve agreement 为什么 collapse，最后把这条教训带回 GLOM 的 local islands。SimCLR 对图像 $x$ 采样两个增强视图 $\tilde x_i,\tilde x_j$，经共享 encoder $f$ 与 projection head $g$ 得到 $\mathbf z_i,\mathbf z_j$。颜色扰动防止模型只用 color histogram 取巧；正样本应接近，不同图像的样本应至少被区分。

\lecturefigure{slide-16-simclr-mechanism.jpg}{SimCLR 流程：两种 crop/颜色增强、共享 encoder、projection 与 agreement。}{Stanford Online 官方视频 00:24:16--00:25:07。}

\begin{knowledgebox}{读图：先分清 $h$ 与 $z$}
左右两路从同一原图采样不同 crop 与颜色扰动，经共享 $f$ 得到 representation $h_i,h_j$，再经 projection head $g$ 得到用于 contrastive objective 的 $z_i,z_j$。图中的“拉近/推远”发生在 projection space；下游 linear probe 常读 encoder representation。这个差别说明训练目标所在空间与最终使用空间可以不同，也提醒我们 GLOM 的 consensus target 必须明确作用在哪个 level 与 state。
\end{knowledgebox}

若只最小化同图 pair 的距离，最容易的解是所有输入都映射到常量 $\mathbf c$：

\[
f(x)=\mathbf c\qquad \forall x.
\]

其中，$f$ 是 encoder，$x$ 是任意输入，$\mathbf c$ 是统一常量表示。这会让所有 positive loss 归零，却完全没有信息。InfoNCE 一类 objective 通过同 batch negatives 或其他 anti-collapse mechanism 阻止该解。这里不展开完整温度归一化推导，重点是“agreement 必须与 separation 同时设计”。

\begin{warningbox}{Agreement 目标的第一风险：全局塌缩}
“让相似位置更相似”不能独立存在。GLOM 需要 local neighborhood、bottom-up/top-down evidence、identity-pose prediction、level separation 与训练 objective 共同阻止整个图、所有层、所有样本变成一个向量。本讲提出了信号来源，但没有给出大规模稳定性证明。
\end{warningbox}

\subsection{Linear probe 只证明可线性读出}

自监督训练后冻结 encoder，再训练一个线性分类器，可以测试表示是否包含容易读出的 label information。好成绩说明 feature 有用，但不说明表示已显式分解 part、whole、pose，也不说明在 generation、robustness 或 causal reasoning 上等价于 supervised model。

\lecturefigure{slide-17-linear-probe-quality.jpg}{用 linear classifier 评估无标签预训练后表示的可读出质量。}{Stanford Online 官方视频 00:25:08--00:28:11。}

\begin{knowledgebox}{Linear probe 的正确读法}
\begin{itemize}
  \item probe 高：类别信息在 frozen representation 中近似线性可分；
  \item probe 低：可能是信息缺失，也可能是信息以非线性方式编码；
  \item probe 无法单独证明：对象边界、part-whole relation、pose equivariance 或 uncertainty calibration。
\end{itemize}
\end{knowledgebox}

\subsection{同图 crop 可能属于不同对象}

如果一张街景同时有汽车、树和行人，把两块 crop 强行拉到同一 vector，可能只学到 scene identity，甚至抹掉 object distinction。GLOM 的目标是：在 scene level 可以大范围一致，在 object level 只让同一 object 的位置一致，在 part level 形成更小 islands。也就是说，agreement 的空间范围必须由 hierarchy 与 evidence 共同决定。

\lecturefigure{slide-18-contrastive-part-mismatch.jpg}{普通 contrastive objective 混淆 scene、object 与 part 的层级。}{Stanford Online 官方视频 00:28:12--00:28:49。}

\lecturefigure{slide-19-spatial-coherence.jpg}{GLOM 希望用 attention 发现不同 coordinate frames 中的局部空间一致性。}{Stanford Online 官方视频 00:28:50--00:30:17。}

\teachervoice{Hinton 把思路追溯到早期 spatial coherence work：相邻输出 vector 在真实表面或深度连续处应当一致。GLOM 的新问题是“不预先知道一致区域在哪里”，所以用 similarity attention 让 region 在 inference 中形成。}

\subsection{本章小结}

Contrastive learning 教会我们：一致性 objective 可以学到强表示，但必须防 collapse，也必须说明“哪一级、哪些位置、在什么条件下应该一致”。GLOM 把这三个问题转交给 hierarchical state、coordinate prediction 与 local attention。

